% from ega1-1-fr.tex
\oldpage[I]{93}
\textit{a}) implique trivialement \textit{b}). Pour voir que \textit{b})
implique \textit{c}), remarquons déjà, puisque $V$ est quasi-compact
(\hyperref[0.2.2.3-fr]{0, 2.2.3}), que $\mathscr{F}$ est isomorphe à un
faisceau $\widetilde{N}|V$, où $N$ est un $A$-module
(\hyperref[I.1.4.1-fr]{1.4.1}). Or, on a
$N=\varinjlim M_\lambda$, où $M_\lambda$ parcourt l'ensemble des
sous-$A$-modules de type fini de $N$, d'où
(\hyperref[I.1.3.9-fr]{1.3.9})
$\mathscr{F}=\widetilde{N}|V=\varinjlim\widetilde{M_\lambda}|V$~; mais comme
$\mathscr{F}$ est de type fini, et $V$ quasi-compact, il existe un indice
$\lambda$ tel que $\mathscr{F}=\widetilde{M_\lambda}|V$
(\hyperref[0.5.2.3-fr]{0, 5.2.3}).

Montrons enfin que \textit{c}) entraîne \textit{a}). Il est clair que
$\mathscr{F}$ est alors de type fini
(\hyperref[I.1.3.6-fr]{1.3.6} et \hyperref[I.1.3.9-fr]{1.3.9})~; en outre,
la question étant locale, on peut se borner au cas où $V=D(f)$, $f\in A$.
Comme $A_f$ est noethérien, on voit finalement que tout revient à prouver
que le noyau d'un homomorphisme $\widetilde{A^n}\to\widetilde{M}$, où $M$
est un $A$-module, est de type fini. Or, un tel homomorphisme est de la
forme $\widetilde{u}$, où $u$ est un homomorphisme $A^n\to M$
(\hyperref[I.1.3.8-fr]{1.3.8}), et si $P=\operatorname{Ker}u$, on a
$\widetilde{P}=\operatorname{Ker}\widetilde{u}$
(\hyperref[I.1.3.9-fr]{1.3.9}). Comme $A$ est noethérien, $P$ est de type
fini, ce qui achève la démonstration.

\begin{corollary}[1.5.2]
\label{I.1.5.2-fr}
Sous les hypothèses de (\hyperref[I.1.5.1-fr]{1.5.1}), le faisceau
$\mathscr{O}_X$ est un faisceau cohérent d'anneaux.
\end{corollary}

\begin{corollary}[1.5.3]
\label{I.1.5.3-fr}
Sous les hypothèses de (\hyperref[I.1.5.1-fr]{1.5.1}), tout faisceau
cohérent sur un ouvert de $X$ est induit par un faisceau cohérent sur $X$.
\end{corollary}

\begin{corollary}[1.5.4]
\label{I.1.5.4-fr}
Sous les hypothèses de (\hyperref[I.1.5.1-fr]{1.5.1}), tout
$\mathscr{O}_X$-Module quasi-cohérent $\mathscr{F}$ est limite inductive
des sous-$\mathscr{O}_X$-Modules cohérents de $\mathscr{F}$.
\end{corollary}

En effet, $\mathscr{F}=\widetilde{M}$ où $M$ est un $A$-module, et $M$ est
limite inductive de ses sous-modules de type fini~; on conclut par
(\hyperref[I.1.3.9-fr]{1.3.9}) et
(\hyperref[I.1.5.1-fr]{1.5.1}).

\subsection{Propriétés fonctorielles des faisceaux quasi-cohérents sur un
spectre premier.}
\label{subsection:I.1.6-fr}

\begin{env}[1.6.1]
\label{I.1.6.1-fr}
Soient $A$, $A'$ deux anneaux,
\[
  \varphi:A'\to A
\]
un homomorphisme,
\[
  {}^a\varphi:X=\operatorname{Spec}(A)\to
  X'=\operatorname{Spec}(A')
\]
l'application continue associée à $\varphi$
(\hyperref[I.1.2.1-fr]{1.2.1}). Nous allons définir un
\emph{homomorphisme canonique}
\[
  \widetilde{\varphi}:\mathscr{O}_{X'}\to
  {}^a\varphi_*(\mathscr{O}_X)
\]
de faisceaux d'anneaux. Pour tout $f'\in A'$, posons $f=\varphi(f')$~; on a
${}^a\varphi^{-1}(D(f'))=D(f)$
(\hyperref[I.1.2.2.2-fr]{1.2.2.2}). Les anneaux
$\Gamma(D(f'),\widetilde{A'})$ et $\Gamma(D(f),\widetilde{A})$
s'identifient respectivement à $A'_{f'}$ et $A_f$
(\hyperref[I.1.3.6-fr]{1.3.6} et \hyperref[I.1.3.7-fr]{1.3.7}). Or,
l'homomorphisme $\varphi$ définit canoniquement un homomorphisme
$\varphi_{f'}:A'_{f'}\to A_f$
(\hyperref[0.1.5.1-fr]{0, 1.5.1}), autrement dit on a un homomorphisme
d'anneaux
\[
  \Gamma(D(f'),\widetilde{A'})\to
  \Gamma({}^a\varphi^{-1}(D(f')),\widetilde{A})
  =\Gamma(D(f'),{}^a\varphi_*(\widetilde{A})).
\]

